\SourceLeaf{267}
\Article{1}
\indent For the problem of
\begin{quote}
mapping the parts of a given surface onto another given surface so that, down to the smallest parts, the image is similar to the original
\end{quote}
\noindent I gave a general solution in 1822, which Conference Councillor \textsc{Schumacher} had printed in issue 3 of the Astronomische Abhandlungen. When this problem is applied to higher geodesy, for which it provides an especially fruitful resource, a need arises for a distinct name for mappings that satisfy the stated condition. I shall therefore call them \textit{conformal} mappings or transfers, giving this otherwise vague adjective a sharply defined mathematical meaning.

In the work cited, the general solution, which contains an arbitrary function, is also applied to several \textit{particular} surfaces. The last example treated there concerns the conformal mapping of the surface of an ellipsoid of revolution onto a spherical surface. It also gives [Art. 13] a choice of the arbitrary function that can be put to very practical use in higher geodesy. That application was only briefly indicated there, with a fuller development left for later. In this memoir, however, I shall develop a slightly modified method for conformally mapping the ellipsoidal surface onto the sphere, much better suited to geodetic applications than \textit{that} particular solution. I shall also include everything required to use it in this way.

\Article{2}
\indent Applying the general solution of the problem to the ellipsoidal and spherical surfaces gives formula (1) below, which covers all conformal mappings from one surface to the other.
\[
T+i\log\operatorname{cotg}\tfrac{1}{2}U=f(t+i\log\{\operatorname{cotg}\tfrac{1}{2}w\mathord{.}(\tfrac{1-e\cos w}{1+e\cos w})^{\frac{1}{2}e}\})
\]

The symbols have the following meanings.
